% ============================================================================
% ARBEITSBLATT DS 1 - Gruppe A: Stadtviertel (S. 30-31)
% ============================================================================

\documentclass[11pt,a4paper]{article}

\usepackage[utf8]{inputenc}
\usepackage[T1]{fontenc}
\usepackage[ngerman]{babel}
\usepackage[left=2cm,right=2cm,top=1.8cm,bottom=1.5cm]{geometry}
\usepackage{helvet}
\usepackage[table]{xcolor}
\usepackage{tabularx}
\usepackage{array}
\usepackage{enumitem}
\usepackage{tcolorbox}
\usepackage{fancyhdr}
\usepackage{lastpage}
\usepackage{amssymb}

\renewcommand{\familydefault}{\sfdefault}

\definecolor{spanisch}{HTML}{c41e3a}
\definecolor{grau}{HTML}{666666}
\definecolor{hellgrau}{HTML}{f5f5f5}

\pagestyle{fancy}
\fancyhf{}
\fancyhead[L]{\small\textcolor{grau}{Spanisch Spätbeginner -- Gruppe A}}
\fancyhead[R]{\small\textcolor{grau}{AB-01-A}}
\fancyfoot[C]{\small\thepage/\pageref{LastPage}}
\renewcommand{\headrulewidth}{0.5pt}

\newcommand{\luecke}[1]{\underline{\hspace{#1}}}
\newcommand{\punktebox}[1]{\hfill\fbox{\small \_\_/#1}}
\newcommand{\es}[1]{\textit{#1}}

\newtcolorbox{merkkasten}[1][]{
    colback=white,
    colframe=spanisch,
    fonttitle=\bfseries\color{spanisch},
    title=#1,
    boxrule=1.5pt,
    arc=0pt,
    left=5pt, right=5pt, top=5pt, bottom=5pt
}

\newtcolorbox{buchverweis}{
    colback=hellgrau,
    colframe=grau,
    boxrule=0.5pt,
    arc=0pt,
    left=5pt, right=5pt, top=3pt, bottom=3pt
}

\newcounter{aufgabe}
\newcommand{\aufgabe}[2]{%
    \stepcounter{aufgabe}%
    \vspace{0.8em}%
    \noindent\textbf{\theaufgabe.} #1 \punktebox{#2}%
    \vspace{0.3em}%
}

\begin{document}

% ============================================================================
% KOPFBEREICH
% ============================================================================
\noindent
\begin{tabularx}{\textwidth}{@{}Xr@{}}
    {\Large\textbf{\textcolor{spanisch}{Mi barrio}}} & Name: \luecke{5cm} \\[0.3em]
    {\small DS 1 -- Bucharbeit (S. 30--31)} & Datum: \luecke{3cm} \\
\end{tabularx}

\vspace{0.3em}
\hrule
\vspace{0.8em}

% ============================================================================
% MERKKASTEN 1: ser / estar / hay
% ============================================================================
\begin{merkkasten}[Kasten 1: ser -- estar -- hay]
\vspace{0.2em}

\begin{tabularx}{\textwidth}{@{}|p{2.5cm}|X|X|@{}}
\hline
\rowcolor{hellgrau}
\textbf{Verb} & \textbf{Verwendung} & \textbf{Beispiel} \\
\hline
\textbf{ser} & \luecke{4cm} (dauerhaft) & Mi barrio \luecke{1.5cm} grande. \\[0.6em]
\hline
\textbf{estar} & \luecke{4cm} (Ort, Zustand) & El cine \luecke{1.5cm} en la calle Mayor. \\[0.6em]
\hline
\textbf{hay} & \luecke{4cm} (es gibt) & \luecke{1.5cm} un parque y muchas tiendas. \\[0.6em]
\hline
\end{tabularx}

\vspace{0.3em}
\end{merkkasten}

\vspace{0.8em}

% ============================================================================
% MERKKASTEN 2: Possessivpronomen vuestro
% ============================================================================
\begin{merkkasten}[Kasten 2: Possessivpronomen \es{vuestro}]
\vspace{0.2em}

\begin{tabularx}{\textwidth}{@{}X|X|X|X@{}}
\textbf{mi} (mein) & \textbf{tu} (dein) & \textbf{su} (sein/ihr/Ihr) & \textbf{nuestro} (unser) \\
\end{tabularx}

\vspace{0.3em}
\hrule
\vspace{0.3em}

\textbf{Neu: vuestro} = euer (Plural, informell -- in Spanien gebräuchlich)

\vspace{0.3em}
\begin{tabularx}{\textwidth}{@{}X|X@{}}
\es{vuestro barrio} (euer Viertel) & \es{vuestra casa} (euer Haus) \\
\es{vuestros padres} (eure Eltern) & \es{vuestras amigas} (eure Freundinnen) \\
\end{tabularx}

\vspace{0.3em}
\end{merkkasten}

\vspace{1em}

% ============================================================================
% AUFGABEN
% ============================================================================

\aufgabe{Ergänze \es{ser}, \es{estar} oder \es{hay}. (Buch S. 30, Aufg. 9)}{6}

\begin{enumerate}[label=\alph*), leftmargin=2em]
    \item Mi barrio \luecke{2cm} muy tranquilo.
    \item El supermercado \luecke{2cm} cerca de mi casa.
    \item En mi barrio \luecke{2cm} dos parques grandes.
    \item La biblioteca \luecke{2cm} moderna y bonita.
    \item ¿Dónde \luecke{2cm} el hospital?
    \item No \luecke{2cm} cine en mi barrio.
\end{enumerate}

\vspace{0.5em}

\aufgabe{Schreibe Sätze mit \es{vuestro/a/os/as}. (Buch S. 30, Aufg. 9)}{4}

\begin{enumerate}[label=\alph*), leftmargin=2em]
    \item (euer Viertel / groß) \luecke{8cm}
    \item (eure Schule / modern) \luecke{8cm}
    \item (eure Freunde / sympathisch) \luecke{8cm}
    \item (eure Häuser / klein) \luecke{8cm}
\end{enumerate}

\vspace{0.5em}

\aufgabe{Wortschatz: Orte im Stadtviertel. Übersetze.}{8}

\begin{tabularx}{\textwidth}{@{}|X|X|X|X|@{}}
\hline
\rowcolor{hellgrau}
\textbf{Spanisch} & \textbf{Deutsch} & \textbf{Spanisch} & \textbf{Deutsch} \\
\hline
el parque & \luecke{2.5cm} & la piscina & \luecke{2.5cm} \\[0.5em]
\hline
el cine & \luecke{2.5cm} & el teatro & \luecke{2.5cm} \\[0.5em]
\hline
la biblioteca & \luecke{2.5cm} & el supermercado & \luecke{2.5cm} \\[0.5em]
\hline
la iglesia & \luecke{2.5cm} & el hospital & \luecke{2.5cm} \\[0.5em]
\hline
\end{tabularx}

\vspace{1em}

\aufgabe{Partnerdialog: Beschreibt euer Stadtviertel. (Buch S. 30, Aufg. 10 + Partnerkarte S. 136)}{6}

\begin{buchverweis}
\textbf{Hinweis:} Nutze die Partnerkarte S. 136. Partner A fragt, Partner B antwortet. Dann Rollentausch.
\end{buchverweis}

\vspace{0.3em}
\textbf{Nützliche Wendungen:}
\begin{itemize}[leftmargin=1.5em, nosep]
    \item \es{¿Cómo es tu barrio?} -- Wie ist dein Viertel?
    \item \es{¿Hay un/una ... en tu barrio?} -- Gibt es ein/eine ... in deinem Viertel?
    \item \es{¿Dónde está el/la ...?} -- Wo ist der/die ...?
\end{itemize}

\vspace{0.3em}
\textbf{Notiere 3 Informationen über das Viertel deines Partners:}

\vspace{0.3em}
1. \luecke{12cm}

\vspace{0.3cm}
2. \luecke{12cm}

\vspace{0.3cm}
3. \luecke{12cm}

\vspace{1em}

\aufgabe{Blogeintrag: Beschreibe dein eigenes Stadtviertel. (Buch S. 31, Aufg. 11 -- Ya sé)}{8}

\begin{buchverweis}
\textbf{Schreibe mindestens 6 Sätze.} Verwende: \es{ser}, \es{estar}, \es{hay}, Orte, Adjektive.
\end{buchverweis}

\vspace{0.5em}
\textbf{Mi barrio}

\vspace{0.3em}
\begin{tabularx}{\textwidth}{@{}|X|@{}}
\hline
\\[0.5em]
\luecke{14cm} \\[0.8em]
\luecke{14cm} \\[0.8em]
\luecke{14cm} \\[0.8em]
\luecke{14cm} \\[0.8em]
\luecke{14cm} \\[0.8em]
\luecke{14cm} \\[0.8em]
\hline
\end{tabularx}

\vspace{1em}
\hrule
\vspace{0.3em}
\begin{flushright}
\textbf{Gesamt: \luecke{1cm} / 32 Punkte}
\end{flushright}

\end{document}
